% !TeX root = ../../Book2.tex
% 纲要标题：### 7.2 内射模
% 纲要位置：计划纲要.md，第 1280 行。

\section{内射模}
\label{b2:sec:0702}

投射模解决“沿满射提升”的问题。若给定的映射定义在一个子模上，希望把它延拓到整个模，所需的是另一种性质。向量空间可以选补空间后延拓线性映射，一般环上的模却未必有补模；内射性把这种延拓能力直接作为研究对象。

% 纲要要点（供目录核对）：
%
% * 延拓性质与 Baer 判据
% * 可除交换群、商群的可除性及内射中间项不保证分裂的例子
% * 模范畴中足够多内射对象
%

\subsection{延拓性质与 Baer 判据}
\label{b2:subsec:0702:01}

\begin{definition}\label{b2:def:injective-module}
设 \(R\) 是环，\(E\) 是左 \(R\) 模。若对任意左 \(R\) 模的子模包含 \(A\subseteq B\) 和同态 \(f:A\to E\)，都存在同态 \(g:B\to E\) 使 \(g|_A=f\)，就称 \(E\) 为\term[neishemo]{内射模}[injective module]。等价地，对任意左 \(R\) 模单同态 \(i:A\to B\) 及 \(f:A\to E\)，都能解出 \(gi=f\)。
\end{definition}

投射性固定同态的来源，沿满射作提升；内射性固定同态的目标 \(E\)，沿单射作延拓。定义中的延拓可画成下面的交换图；虚线是所求的同态，交换性即 \(gi=f\)。
\[
\begin{tikzcd}[column sep=3em,row sep=2.5em]
A\arrow[r,hook,"i"]\arrow[d,"f"']&B\arrow[dl,dashed,"g"]\\
E&
\end{tikzcd}
\]

\(\mathbb Z\) 不是内射 \(\mathbb Z\) 模：把 \(2\mathbb Z\) 映到 \(\mathbb Z\) 的同态 \(2n\mapsto n\) 若能延拓，则 \(2g(1)=1\)，在整数中无解。定义中的“内射”描述目标模的延拓性质，并不表示该模中的所有同态都是单射。

\begin{proposition}\label{b2:prop:injective-characterizations}
设 \(R\) 是环，\(E\) 是左 \(R\) 模。\(E\) 内射，等价于 \(\operatorname{Hom}_R(-,E)\) 将每个左 \(R\) 模的短正合列反向变成交换群的短正合列；也等价于每个左 \(R\) 模短正合列 \(0\to E\to B\to C\to0\) 都分裂。
\end{proposition}
\begin{proof}
\cref{b2:prop:hom-left-exact}说明 Hom 列只可能在最后一项不满，满射性恰是子模上的同态均能延拓。内射时，将恒等映射 \(E\to E\) 沿 \(E\to B\) 延拓，得到缩回，原列由\cref{b2:thm:split-exact-sequence}分裂。

反过来，设所有左端为 \(E\) 的短正合列分裂。给定左 \(R\) 模的子模包含 \(A\subseteq B\) 及同态 \(f:A\to E\)，希望把 \(E\) 和 \(B\) 放进同一个模，并强制 \(B\) 中的 \(a\) 与 \(E\) 中的 \(f(a)\) 相等。若 \(E\) 在这个模中有缩回，将 \(B\) 中的元素送回 \(E\) 就可得到延拓。为此，在 \(E\oplus B\) 中加入关系 \((f(a),-a)=0\)，按\cref{b2:prop:module-pullback-pushout}的推出构造取商模
\[
D=(E\oplus B)/\Bigl\{\,\bigl(f(a),-a\bigr):a\in A\,\Bigr\}.
\]
分母集合是同态 \(a\mapsto\bigl(f(a),-a\bigr)\) 的像，因而是子模。记商映射为 \(\pi:E\oplus B\to D\)。映射 \(j:E\to D\)、\(e\mapsto\pi(e,0)\) 单射：若 \((e,0)=\bigl(f(a),-a\bigr)\)，则 \(a=0,e=0\)。故 \(0\to E\xrightarrow{j}D\to D/j(E)\to0\) 有缩回 \(r:D\to E\)。令 \(g(b)=r\bigl(\pi(0,b)\bigr)\)，则对 \(a\in A\)，商中的等式 \(\pi(0,a)=\pi\bigl(f(a),0\bigr)\) 给出 \(g(a)=f(a)\)。
\end{proof}

这一构造把任意延拓问题改写为一个以 \(E\) 开头的分裂问题。与\cref{b2:thm:projective-characterizations}相比，投射性使所有以该模为右端的短正合列分裂，内射性则使所有以该模为左端的短正合列分裂。判断内射性时，还可以把需要检验的子模缩小到环自身的左理想：一次加入一个新元素时，它与旧定义域之间的全部关系恰由一个左理想记录。

\begin{theorem}[Baer 判据]\label{b2:thm:baer-criterion}
设 \(R\) 是含幺环，\(E\) 是左 \(R\) 模。\(E\) 内射，当且仅当对每个左理想 \(I\subseteq R\)，每个左 \(R\) 模同态 \(f:I\to E\) 都能延拓为左 \(R\) 模同态 \({}_RR\to E\)。后一条件也可写成：对每个这样的 \(I,f\)，存在 \(e\in E\)，使所有 \(a\in I\) 都满足 \(f(a)=ae\)。
\end{theorem}

这称为\term[baer-panju]{Baer 判据}[Baer's criterion]。等式形式来自 \(R\to E\) 的同态由 \(1\) 的像唯一确定。注意这里检验所有左理想；除非环有额外有限性条件，不能只检验有限生成的左理想。

\begin{proof}
必要性直接取 \(A=I,B=R\)。为证充分性，给定左 \(R\) 模的子模包含 \(A\subseteq B\) 和同态 \(f:A\to E\)，考虑所有延拓对 \((N,h)\)，其中 \(A\subseteq N\subseteq B\)、\(h:N\to E\)、\(h|_A=f\)。以定义域包含且映射限制相容为序。原始对 \((A,f)\) 保证集合非空；一条链的定义域之并仍是子模，相容的映射在其上拼成同态，给出上界；空链也有 \((A,f)\) 作上界。用第一卷附录A定理A.1.4的 Zorn 引理，取极大对 \((N,h)\)。

若 \(N\ne B\)，取 \(b\in B\setminus N\)，准备把 \(h\) 延拓到 \(N+Rb\)。若指定 \(h'(b)=e\)，则凡是 \(ab\in N\)，新旧映射都必须满足 \(h(ab)=ae\)。这些使新元素落回旧定义域的标量及旧映射强迫它们满足的条件，给出
\[
I=\{\,a\in R:ab\in N\,\},\qquad u:I\longrightarrow E,\quad u(a)=h(ab).
\]
\(I\) 对加法封闭，且 \(a\in I,r\in R\) 时，\((ra)b=r(ab)\in N\)，所以 \(I\) 是左理想；\(h\) 的线性又使 \(u\) 成为左模同态。按假设取 \(e\in E\) 使 \(u(a)=ae\)。一个延拓若把 \(b\) 送到 \(e\)，线性就强迫它只能按下式取值：
\[
h'(n+rb)=h(n)+re.
\]
所剩问题是这个公式会不会依赖 \(n+rb\) 的表示。若 \(n+rb=n'+r'b\)，则 \(a=r-r'\in I\) 且 \(ab=n'-n\)，所以 \(h(n'-n)=ae=(r-r')e\)，即
\[
h(n)+re=h(n')+r'e.
\]
因此 \(h'\) 良定义。加法与左标量相容由公式直接验证，且 \(h'|_N=h\)。它把极大延拓的定义域严格扩大，矛盾。因此 \(N=B\)，得到所需延拓。
\end{proof}

\subsection{可除交换群}
\label{b2:subsec:0702:02}

\begin{definition}\label{b2:def:divisible-group}
设 \(D\) 是交换群。若对每个整数 \(n>0\) 和每个 \(d\in D\)，都存在 \(x\in D\) 使 \(nx=d\)，就称 \(D\) 为\term[kechu-qun]{可除群}[divisible group]。这要求乘 \(n\) 为满射，不要求它为单射。
\end{definition}

\(\mathbb Q\) 与 \(\mathbb Q/\mathbb Z\) 都可除：分别用有理数除以 \(n\)，或把有理数代表元除以 \(n\)。后者有大量有限阶元素，说明可除与无挠是不同条件。非零有限交换群不可能可除，因为取群阶 \(n\)，乘 \(n\) 为零映射。

\ProseParagraphBreak
在整数环上，Baer 判据的延拓问题恰好变为这些除法方程。每个非零理想都是 \(n\mathbb Z\)，其中 \(n>0\)；给定 \(f:n\mathbb Z\to D\)，一个延拓由 \(x=g(1)\) 决定，所需条件只有 \(nx=f(n)\)。

\begin{theorem}\label{b2:thm:divisible-injective}
设 \(D\) 是交换群，并按整数倍看成左 \(\mathbb Z\) 模，则 \(D\) 内射，当且仅当它可除。若 \(D\) 可除，则对每个子群 \(A\leq D\)，商群 \(D/A\) 也可除，因而内射。但是交换群的短正合列
\[
0\longrightarrow\mathbb Z\longrightarrow\mathbb Q
\longrightarrow\mathbb Q/\mathbb Z\longrightarrow0,
\]
其中箭头分别为包含及自然商映射，并不分裂；它的中间项和右端均为内射 \(\mathbb Z\) 模。
\end{theorem}
\begin{proof}
若 \(D\) 内射，给定 \(d\in D\) 与 \(n>0\)，同态 \(n\mathbb Z\to D\)、\(nk\mapsto kd\) 可延拓为 \(g:\mathbb Z\to D\)。取 \(x=g(1)\)，便有 \(nx=g(n)=d\)。

若 \(D\) 可除，用\cref{b2:thm:baer-criterion}检验整数理想。零理想的映射用零映射延拓；对 \(n\mathbb Z\)、\(n>0\)，给定 \(f\)，取 \(x\) 使 \(nx=f(n)\)。令 \(g(k)=kx\)，则 \(g(nk)=k\cdot f(n)=f(nk)\)，得到延拓。

再设 \(D\) 可除，\(A\leq D\)。给定 \(d+A\in D/A\) 和正整数 \(n\)，由可除性取 \(x\in D\) 使 \(nx=d\)，则 \(n(x+A)=d+A\)。所以商群可除，由刚证得的等价性又知它内射。特别地，\(\mathbb Q\) 可除，其商 \(\mathbb Q/\mathbb Z\) 也可除，两者都内射。

若所写短正合列分裂，\cref{b2:thm:split-exact-sequence}给出缩回 \(\rho:\mathbb Q\to\mathbb Z\)，且 \(\rho(1)=1\)。但群同态又要求 \(2\rho(1/2)=\rho(1)=1\)，在整数中不可能。因此该列不分裂。这不违背内射模的分裂刻画：它的左端是并不内射的 \(\mathbb Z\)。
\end{proof}

\(\mathbb Z\) 投射但不内射。在域上，每个向量空间都有基和补空间，所以每个模既投射又内射；离开域以后，这两个条件分别约束同态的出发端和到达端。

\subsection{模范畴中足够多内射对象}
\label{b2:subsec:0702:03}

每个左 \(R\) 模 \(M\) 都有自由模的满射 \(F\twoheadrightarrow M\)，自由模又投射，所以模范畴有足够多投射对象。内射方向需要证明相反的存在性：总能找到内射模 \(E\) 及单射 \(M\hookrightarrow E\)。

\ProseParagraphBreak
先从底交换群着手。若要用映入某个交换群的同态检测非零元素，\(\mathbb Q\) 无法检测有限阶元素；\(\mathbb Q/\mathbb Z\) 却含有每种有限阶所需的像，又是内射交换群，可以把循环子群上的同态延拓到整个群。这给出下述分离性质。

\begin{lemma}\label{b2:lem:qz-separates}
设 \(A\) 是交换群，\(0\ne a\in A\)。存在群同态 \(\chi:A\to\mathbb Q/\mathbb Z\) 使 \(\chi(a)\ne0\)。
\end{lemma}
\begin{proof}
先在 \(a\) 生成的循环子群上定义映射。若 \(a\) 的阶为 \(n>1\)，令 \(a\mapsto 1/n+\mathbb Z\)；若 \(a\) 为无限阶，令 \(a\mapsto1/2+\mathbb Z\)。两种情况均给出循环群同态，且 \(a\) 的像非零。\cref{b2:thm:divisible-injective}说明 \(\mathbb Q/\mathbb Z\) 内射，所以该同态可延拓到 \(A\)。
\end{proof}

然后把交换群上的延拓转回 \(R\) 模上的延拓。对保单位元的同态 \(\mathbb Z\to R\)，限制标量就是取底交换群；\cref{b2:prop:scalar-change-adjunctions}给出的右伴随是余诱导 \(D\mapsto\operatorname{Hom}_{\mathbb Z}(R,D)\)。它将映入 \(D\) 的群同态与映入余诱导模的 \(R\) 线性映射相对应，因此可用 \(D\) 的内射性构造一个内射 \(R\) 模。

\begin{lemma}\label{b2:lem:coinduced-injective}
设 \(R\) 是含幺环，\(D\) 是作为 \(\mathbb Z\) 模内射的交换群，则
\[
J=\operatorname{Hom}_{\mathbb Z}(R,D),\qquad (rf)(s)=f(sr)
\]
是内射左 \(R\) 模。
\end{lemma}
\begin{proof}
这里 \(R\) 在 Hom 中先作为加法群。由\cref{b2:prop:scalar-change-adjunctions}，所写作用使 \(J\) 成为左 \(R\) 模；对任意左 \(R\) 模 \(M\)，其伴随同构给出互逆对应
\[
\begin{aligned}
\operatorname{Hom}_R(M,J)&\longleftrightarrow\operatorname{Hom}_{\mathbb Z}(M,D),\\
\alpha&\longmapsto\bigl(m\mapsto\alpha(m)(1)\bigr),\\
\chi&\longmapsto\Bigl[m\mapsto\bigl(s\mapsto\chi(sm)\bigr)\Bigr].
\end{aligned}
\]
对给定的 \(\chi\)，记第二个公式所得的左 \(R\) 线性映射为 \(\alpha_\chi\)，即 \(\alpha_\chi(m)(s)=\chi(sm)\)。

若 \(A\subseteq B\) 是左 \(R\) 模的子模包含，给定 \(A\to J\)，把它变为 \(A\to D\) 的加法群同态 \(\chi\)。由 \(D\) 内射，延拓为 \(\widetilde\chi:B\to D\)，再用反向公式得到 \(\alpha_{\widetilde\chi}:B\to J\)。在 \(a\in A,s\in R\) 处有 \(\widetilde\chi(sa)=\chi(sa)\)，所以后者确实延拓原映射。
\end{proof}

\begin{theorem}[足够多内射对象]\label{b2:thm:enough-injectives}
对任意含幺环 \(R\)，每个左 \(R\) 模 \(M\) 都能嵌入一个内射左模。
\end{theorem}
\begin{proof}
令 \(D=\mathbb Q/\mathbb Z\)、\(J=\operatorname{Hom}_{\mathbb Z}(R,D)\)。先对每个底交换群同态 \(\chi:M\to D\)，按\cref{b2:lem:coinduced-injective}中的伴随对应取左 \(R\) 线性映射 \(\alpha_\chi:M\to J\)，其中 \(\alpha_\chi(m)(s)=\chi(sm)\)。一个这样的映射未必检测所有非零元素；于是以集合 \(H=\operatorname{Hom}_{\mathbb Z}(M,D)\) 为指标，将它们全部组成积映射
\[
\iota:M\longrightarrow\prod_{\chi\in H}J,
\qquad \iota(m)=\bigl(\alpha_\chi(m)\bigr)_{\chi\in H}.
\]
积的泛性质保证 \(\iota\) 左线性。若 \(m\ne0\)，\cref{b2:lem:qz-separates}给出 \(\chi(m)\ne0\) 的分量；在 \(s=1\) 处有 \(\alpha_\chi(m)(1)=\chi(m)\ne0\)，故 \(\iota\) 单射。

由\cref{b2:thm:divisible-injective}及\cref{b2:lem:coinduced-injective}，每个 \(J\) 内射，任意积 \(\prod J\) 也内射：对子模上的映射逐分量延拓，再用\cref{b2:thm:module-product-universal}合成一个映入积的同态。这种合成不要求有限支集；分量族无限时，选择各延拓使用选择公理。因此所构造的积就是所需内射模。
\end{proof}

内射模的积封闭性与\cref{b2:prop:projective-sums-summands}中的投射模直和封闭性相对应：映入积的同态按坐标延拓，从直和出发的同态按分量提升，分别使用积与余积的泛性质。

\ProseParagraphBreak
“\term[zugouduo-neisheduixiang]{足够多内射对象}[enough injectives]”意指上述嵌入总能找到；不意指每个模本身内射，也不意指内射模可以选得有限生成。例如，有限生成可除交换群只能为零：按\cref{b2:thm:pid-module-structure}分成有限秩自由部分与有限挠部分后，乘二的满射性排除非零自由部分，再乘挠部分的阶便排除非零挠部分。因此任何包含 \(\mathbb Z\) 的内射交换群都不可能有限生成。反复对商模作这样的嵌入，将在第四卷构造内射消解。
